% ============================================
% 第3章：参考文献/图片/文档导出模块概要设计
% 负责人：C
% ============================================

\section{参考文献管理模块概要设计}

\subsection{模块职责}

参考文献管理模块负责学生论文写作过程中参考文献条目的全生命周期管理，包括文献信息的CRUD操作、分类检索和标准格式化输出。该模块独立于论文实体运行，维护一个通用的参考文献库，导出时按序注入文档。

\subsection{模块接口}

\begin{table}[htbp]
\centering
\caption{参考文献管理模块API接口}
\label{tab:ref_api}
\begin{tabular}{lllp{5cm}}
\toprule
\textbf{方法} & \textbf{路径} & \textbf{说明} \\
\midrule
GET    & /api/references       & 获取所有文献（按年份降序） \\
GET    & /api/references/\{id\} & 获取单条文献 \\
POST   & /api/references       & 新增文献 \\
PUT    & /api/references/\{id\} & 更新文献 \\
DELETE & /api/references/\{id\} & 删除文献 \\
GET    & /api/references/search/author & 按作者搜索 \\
GET    & /api/references/search/title  & 按标题搜索 \\
GET    & /api/references/\{id\}/format & 单条GB/T 7714格式化 \\
GET    & /api/references/format/all   & 全部GB/T 7714格式化列表 \\
\bottomrule
\end{tabular}
\end{table}

\subsection{与外部模块的交互}

\begin{itemize}
    \item \textbf{与论文管理模块}：论文的参考文献数据通过thesis\_id外键关联，导出模块读取格式化后的参考文献列表写入文档；
    \item \textbf{与前端交互}：前端通过REST API完成文献的增删改查操作，格式化结果在前端预览列表展示；
    \item \textbf{与导出模块}：导出模块调用格式化服务获取GB/T 7714格式化的文献字符串列表，按序写入DOCX文档的"参考文献"章节。
\end{itemize}

\subsection{GB/T 7714格式化设计}

系统支持五种文献类型的标准化格式化输出：

\begin{itemize}
    \item \textbf{[J] 期刊文章}：作者. 标题[J]. 刊物名称, 年份, 卷(期): 页码.
    \item \textbf{[M] 专著/书籍}：作者. 标题[M]. 出版地: 出版社, 年份.
    \item \textbf{[C] 会议论文}：作者. 标题[C]// 会议名称. 年份: 页码.
    \item \textbf{[D] 学位论文}：作者. 标题[D]. 授予单位, 年份.
    \item \textbf{[EB] 网络资源}：作者. 标题[EB/OL]. 链接, 引用日期.
\end{itemize}

格式化器采用策略模式实现，根据文献类型自动路由到对应的格式化方法，便于后续扩展新的文献类型。

\section{图片上传与管理模块概要设计}

\subsection{模块职责}

图片管理模块负责学生论文插图的文件上传、校验、存储和访问。文件存储于服务器本地文件系统，图片元信息存储于数据库。

\subsection{模块接口}

\begin{table}[htbp]
\centering
\caption{图片管理模块API接口}
\label{tab:img_api}
\begin{tabular}{lllp{5cm}}
\toprule
\textbf{方法} & \textbf{路径} & \textbf{说明} \\
\midrule
POST   & /api/images/upload       & 上传图片（multipart/form-data） \\
GET    & /api/images/\{id\}        & 获取图片元信息 \\
GET    & /api/images/\{id\}/file   & 获取图片文件（浏览器直接查看） \\
DELETE & /api/images/\{id\}        & 删除图片 \\
\bottomrule
\end{tabular}
\end{table}

\subsection{技术方案}

\begin{itemize}
    \item \textbf{上传协议}：HTTP multipart/form-data，通过Spring MVC的MultipartFile接收；
    \item \textbf{文件存储}：服务器本地文件系统，存储路径为 \verb|{uploadDir}/{uuid}.{ext}|；
    \item \textbf{文件名策略}：采用UUID v4生成存储文件名，避免中文文件名编码问题和重名冲突；
    \item \textbf{格式校验}：后端校验Content-Type头，仅允许image/jpeg、image/png、image/gif、image/webp；
    \item \textbf{大小限制}：通过spring.servlet.multipart.max-file-size=10MB全局限制；
    \item \textbf{访问方式}：通过/api/images/\{id\}/file静态资源映射返回文件流，Content-Disposition设置为inline以支持浏览器直接展示。
\end{itemize}

\subsection{与外部模块的交互}

\begin{itemize}
    \item \textbf{与前端交互}：前端在上传成功后获取图片URL，将其嵌入编辑器HTML内容中；
    \item \textbf{与导出模块}：导出模块解析论文章节HTML，提取<img>标签中的图片URL，读取本地图片文件嵌入DOCX/PDF文档。
\end{itemize}

\section{文档导出模块概要设计}

\subsection{模块职责}

文档导出模块是本系统的核心模块，负责将学生在Web编辑器中撰写的论文内容，按模板格式配置生成标准的Word（DOCX）和PDF文档。

\subsection{总体流程}

\begin{enumerate}
    \item 学生点击"导出"按钮，前端发送导出请求（指定论文ID和导出格式）；
    \item 后端接收请求，校验论文状态和内容完整性；
    \item 后端从数据库读取论文数据（各章节内容、参考文献列表、模板格式配置）；
    \item 获取嵌入的图片文件资源（从本地文件系统读取）；
    \item 使用Apache POI XWPF按模板格式配置创建DOCX文档；
    \item 若导出格式为PDF，则调用LibreOffice无头模式将DOCX转换为PDF；
    \item 将生成的文档文件返回给前端（直接下载或返回文件链接）。
\end{enumerate}

\subsection{技术方案对比}

\begin{table}[htbp]
\centering
\caption{文档生成技术方案对比}
\label{tab:export_compare}
\begin{tabular}{|c|c|c|c|}
\hline
\textbf{方案} & \textbf{DOCX生成} & \textbf{PDF生成} & \textbf{中文支持} \\
\hline
Apache POI & 原生创建、样式精细 & 不支持 & 需指定中文字体 \\
\hline
LibreOffice Headless & — & 从DOCX转换，还原度高 & 系统字体依赖 \\
\hline
\end{tabular}
\end{table}

技术选型结论：
\begin{itemize}
    \item DOCX生成采用Apache POI（poi-ooxml 5.x），直接操作XWPFDocument对象，逐段落、表格、图片构建文档；
    \item PDF生成采用"DOCX → LibreOffice Headless → PDF"的两步方案，保证DOCX和PDF输出的一致性。
\end{itemize}

\subsection{数据流设计}

\begin{verbatim}
用户点击导出
    |
    v
导出请求 -> 查询论文数据（Thesis + ThesisSections）
    |              |
    |       查询参考文献（ThesisReference）
    |              |
    |       查询图片文件（ThesisImage → 本地文件）
    |              |
    v              v
Apache POI XWPF 构建 .docx
    |
    +---> 直接返回 .docx（格式：DOCX）
    |
    +---> LibreOffice --headless --convert-to pdf
              |
              v
          返回 .pdf（格式：PDF）
\end{verbatim}

\subsection{与外部模块的交互}

\begin{itemize}
    \item \textbf{与论文管理模块}：读取论文的章节内容（HTML格式）和模板格式参数（JSON）；
    \item \textbf{与参考文献模块}：调用格式化服务，获取GB/T 7714格式化的参考文献字符串列表；
    \item \textbf{与图片管理模块}：通过图片存储路径读取本地图片文件，嵌入文档对应位置；
    \item \textbf{与前端交互}：通过异步任务接口反馈导出进度，最终返回下载链接。
\end{itemize}
